\documentclass[a4paper,12pt]{article}
\usepackage{_document_style_article_eng}


\title{%
  \textbf{Barbalat's Lemma}\\
  %\rule[-1cm]{0pt}{1cm} % linea vuota per spaziatura
}
\author{%
  Davide De Marco
  %\\{\relsize{-1}Bacoli, Italia}
}
\date{}

\begin{document}
\maketitle
\section{Statement and Proof of Barbalat's Lemma}
\begin{namelemma}[Barbalat's]\label{lemma:barbalat}
  Sia $f \, \colon \interval[open right]{0}{+\infty} \to \mathbb{R}$ una funzione tale che:
  \begin{enumerate}[label=\upshape({\itshape\roman*}\kern 0.6pt), ref=({\itshape\roman*}\kern 0.6pt)]
    \item $f \in \classC{1}\bigl(\interval[open right]{0}{+\infty}\bigr)$;
    \item $\lim\limits_{t \to + \infty} f(t)$ esiste e vale $\mathscr{l} \in \mathbb{R}$;
    \item $\dot{f}$ è uniformemente continua su $\interval[open right]{0}{+\infty}$.
  \end{enumerate}
  Allora, $\lim\limits_{t \to + \infty} \dot{f}(t) = 0$.
\end{namelemma}

\begin{proof}
  Si supponga per assurdo che 
  \[
  \lim_{t \to +\infty} \dot{f}(t) \neq 0
  \]
  oppure che $\dot{f}$ non ammetta limite per $t \to +\infty$. In entrambi i casi, negando la definizione di limite, esiste una costante $\varepsilon > 0$ e una successione $\mleft\{ t_n \mright\}_{n \in \mathbb{N}}$ tale che:
  \begin{equation}\label{eq:barbalat_contr}
    \lim_{n \to +\infty}t_n = +\infty \quad \text{e} \quad \mleft|\dot{f}(t_n) \mright| > \varepsilon \quad \forall n \in \mathbb{N}.
  \end{equation}
  Poiché $\dot{f}$ è uniformemente continua, per ogni $\bar{\varepsilon} > 0$ esiste $\delta = \delta(\bar{\varepsilon}) > 0$ tale che
  \begin{equation*}
    \mleft| \dot{f}(t) - \dot{f}(t_n) \mright| < \bar{\varepsilon} \quad \forall t \in \interval[open]{t_n}{t_n + \delta}, \,\forall n \in \mathbb{N}.
  \end{equation*}
  Scegliendo $\bar{\varepsilon} = \varepsilon/2$ si ottiene la relazione
  \begin{equation}\label{eq:barbalat_uc2}
    \mleft| \dot{f}(t) - \dot{f}(t_n) \mright| < \frac{\varepsilon}{2} \quad \forall t \in \interval[open]{t_n}{t_n + \delta}, \,\forall n \in \mathbb{N}.
  \end{equation}
  Mediante la disuguaglianza triangolare si ricava la stima
  \begin{equation*}
    \mleft| \dot{f}(t_n) \mright| = \mleft|\dot{f}(t) - \dot{f}(t_n) - \dot{f}(t) \mright| \leq \mleft| \dot{f}(t) - \dot{f}(t_n) \mright| + \mleft| \dot{f}(t) \mright|
  \end{equation*}
  da cui segue necessariamente
  \begin{equation}\label{eq:barbalat_triang2}
    \mleft| \dot{f}(t) \mright| \geq \mleft| \dot{f}(t_n) \mright| - \mleft| \dot{f}(t) - \dot{f}(t_n) \mright| \quad \forall t \in \interval[open]{t_n}{t_n + \delta}, \,\forall n \in \mathbb{N}.
  \end{equation}
  Combinando \eqref{eq:barbalat_contr}, \eqref{eq:barbalat_uc2} e \eqref{eq:barbalat_triang2} si ottiene
  \begin{equation*}
    \mleft| \dot{f}(t) \mright| \geq \mleft| \dot{f}(t_n) \mright| - \mleft| \dot{f}(t) - \dot{f}(t_n) \mright| > \varepsilon - \frac{\varepsilon}{2} = \frac{\varepsilon}{2}
  \end{equation*}
  da cui segue che
  \begin{equation}\label{eq:barbalat_contr3}
    \mleft| \dot{f}(t) \mright| > \frac{\varepsilon}{2} \quad \forall t \in \interval[open]{t_n}{t_n + \delta}, \,\forall n \in \mathbb{N}.
  \end{equation}
  Per ipotesi, per ogni $n \in \mathbb{N}$, $f$ è continua su $\interval{t_n}{t_n + \delta}$ e derivabile su $\interval[open]{t_n}{t_n + \delta}$, dunque, per il \thmcita{Teorema di Lagrange} e la relazione \eqref{eq:barbalat_contr3}, esiste $\xi_n \in \interval[open]{t_n}{t_n + \delta}$ tale che
  \begin{equation*}
    \mleft|f(t_n + \delta) - f(t_n)\mright| = \mleft|\dot{f}(\xi_n) \cdot \delta \mright| = \mleft|\dot{f}(\xi_n)\mright| \delta > \frac{\varepsilon\delta}{2}
  \end{equation*}
  da cui segue che
  \begin{equation}\label{eq:barbalat_contr4}
    \mleft|f(t_n + \delta) - f(t_n)\mright| > \frac{\varepsilon\delta}{2} > 0 \quad \forall n \in \mathbb{N}.
  \end{equation}
  Per ipotesi
  \[
  \lim_{n \to +\infty} f(t_n + \delta) - f(t_n) = \mathscr{l} - \mathscr{l} = 0,
  \]
  applicando la definizione di limite:
  \begin{equation*}
    \forall \sigma > 0, \, \exists \nu \in \mathbb{N} \, \colon \, \mleft|f(t_n + \delta) - f(t_n)\mright| < \sigma \quad \forall n > \nu. 
  \end{equation*}
  Scegliendo $\sigma = \varepsilon\delta/2$ si ottiene la relazione
  \begin{equation*}
    \mleft|f(t_n + \delta) - f(t_n)\mright| < \frac{\varepsilon\delta}{2} \quad \forall n > \nu
  \end{equation*}
  che è in contraddizione con \eqref{eq:barbalat_contr4}. L'assurdo nasce dall'aver negato la tesi, dunque, necessariamente, 
  \[
  \lim_{t \to + \infty} \dot{f}(t) = 0.
  \]
\end{proof}

\medskip
Se la funzione $\dot{f}$ non è uniformemente continua, non c'è garanzia che la tesi del \thmcita{Lemma di Barbalat \ref{lemma:barbalat}} sia vera. Viene mostrato, di seguito, un controesempio a sostegno di questa affermazione.

\section{Counterexample to the Barbalat's Lemma}
\begin{example}
Si consideri la funzione
\begin{equation}
    f(t) = \frac{\sin t^2}{t}
\end{equation}
definita su $\interval[open right]{0}{+\infty}$. Per il \thmcita{Teorema dei Carabinieri} si osserva subito che
\begin{equation*}
    \lim_{t \to +\infty} f(t) = 0 \,.
\end{equation*}
Calcolando la derivata prima si ottiene:
\begin{equation*}
    \dot{f}(t) = \diff*{\mleft\{\frac{\sin t^2}{t}\mright\}}{t} = \frac{2t \cos t^2 \cdot t - \sin t^2}{t^2} = 2\cos t^2 - \frac{\sin t^2}{t^2}
\end{equation*}

\smallskip
Siano definite le seguenti successioni:
\begin{equation*}
    x_n = \SQRT{2n\pi}, \quad y_n = \SQRT{2n\pi + \frac{\pi}{2}} \,.
\end{equation*}
La successione definita come $x_n - y_n$ è della forma
\begin{equation*}
    x_n - y_n = \SQRT{2n\pi} - \SQRT{2n\pi + \frac{\pi}{2}} = \frac{-\frac{\pi}{2}}{\SQRT{2n\pi} + \SQRT{2n\pi + \frac{\pi}{2}}} \,,
\end{equation*}
da cui segue che
\begin{equation}\label{eq:barbalat_contr5}
    \lim_{n \to +\infty} \mleft| x_n - y_n \mright| = 0.
\end{equation}
Si osserva che, per ogni $n \in \mathbb{N}$, $\dot{f}$ valutata in $x_n$ vale
\begin{equation*}
    \dot{f}(x_n) = 2\cos(2n\pi) - \frac{\sin(2n\pi)}{2n\pi} = 2(1) - 0 = 2 \, ,
\end{equation*}
mentre, per ogni $n \in \mathbb{N}$, $\dot{f}$ valutata in $y_n$ vale
\begin{equation*}
    \dot{f}(y_n) = 2\cos\mleft(2n\pi + \frac{\pi}{2}\mright) - \frac{\sin\mleft(2n\pi + \frac{\pi}{2}\mright)}{2n\pi + \frac{\pi}{2}} = 2(0) - \frac{1}{2n\pi + \frac{\pi}{2}} = - \frac{1}{2n\pi + \frac{\pi}{2}} \,.
\end{equation*}
Dunque, 
\begin{equation*}
  \mleft| \dot{f}(x_n) - \dot{f}(y_n) \mright| > 2 \quad \forall n \in \mathbb{N}
\end{equation*}
e, in particolare,
\begin{equation*}
    \lim_{n \to +\infty} \mleft| \dot{f}(x_n) - \dot{f}(y_n) \mright| = \lim_{n \to +\infty} \mleft| 2 + \frac{1}{2n\pi + \frac{\pi}{2}} \mright| = 2 \,.
\end{equation*}

Si scelga $\varepsilon = 2 > 0$ e si fissi un arbitrario $\delta > 0$. 
Dalla \eqref{eq:barbalat_contr5} segue che
\begin{equation*}
    \exists \nu \in \mathbb{N} \, \colon \, \mleft| x_n - y_n \mright| < \delta \quad \forall n > \nu.
\end{equation*}
Fissando un qualsiasi indice $\bar{n} > \nu$, si ottiene la relazione
\begin{equation*}
    \mleft| x_{\bar{n}} - y_{\bar{n}} \mright| < \delta.
\end{equation*}
Tuttavia, valutando la distanza tra le rispettive immagini, si osserva che
\begin{equation*}
    \mleft| \dot{f}(x_{\bar{n}}) - \dot{f}(y_{\bar{n}}) \mright| = 2 + \frac{1}{2\bar{n}\pi + \frac{\pi}{2}} > 2 = \varepsilon \,.
\end{equation*}
Si è così esibita una coppia di punti che, pur distando meno dell'arbitraria tolleranza $\delta$, possiedono immagini che distano strettamente più di $\varepsilon$. Essendo ciò costruttivamente valido per ogni possibile scelta di $\delta > 0$, la definizione di uniforme continuità per $\dot{f}$ risulta formalmente violata.

\smallskip
Infine, le medesime successioni $x_n$ e $y_n$ permettono di dimostrare in modo immediato che $\dot{f}$ non ammette limite per $t \to +\infty$. Si osserva banalmente che
\begin{equation*}
    \lim_{n \to +\infty} x_n = +\infty \quad \text{e} \quad \lim_{n \to +\infty} y_n = +\infty \,,
\end{equation*}
ma, passando al limite le rispettive immagini, si ottiene:
\begin{align*}
    \lim_{n \to +\infty} \dot{f}(x_n) &= \lim_{n \to +\infty} 2 = 2 \\
    \lim_{n \to +\infty} \dot{f}(y_n) &= \lim_{n \to +\infty} \mleft( - \frac{1}{2n\pi + \frac{\pi}{2}} \mright) = 0 \,
\end{align*}
da cui si deduce necessariamente che
\begin{equation*}
    \nexists \lim_{t \to +\infty} \dot{f}(t) \,.
\end{equation*}
\end{example}
\end{document}